\textbf{\textit{Soal. }}Diberikan segitiga lancip $ABC$. Garis $\ell_1, \ell_2 \perp AB$ berturut-turut pada $A$ dan $B$. Misalkan $M$ titik tengah $AB$. Garis melalui $M$ tegak lurus $AC$ dan $BC$ memotong $\ell_1$ dan $\ell_2$ berturut-turut pada $E$ dan $F$. Misalkan $D$ titik potong dari $EF$ dan $MC$. Jika $EF = \frac{64}{3}$, $FM = 16$, dan $BD = 12$. Luas segitiga $AMF$ dapat dinyatakan dalam $a\sqrt{b}$ dimana $a$ dan $b$ bilangan asli dan $b$ tidak dapat dibagi oleh bilangan kuadrat sempurna selain 1. Nilai dari $a + b$ adalah \dots
\begin{proof}[\textbf{Solusi. }] (Senin, 13 November 2023) Misalkan $G$ dan $H$ berturut-turut adalah perpotongan $EM$ dengan $AC$ dan $MF$ dengan $CB$. 
\begin{center}
    \begin{asy}
import olympiad;
import geometry;
unitsize(1cm);
pair A, B, C, M, D, E, F, G, H, la, lb;
A = (-3,0);
B = (3,0);
la = (-3,5);
lb = (3, 5);
M = (A+B)/2;
C = (-1,5);
G = foot(M, A, C);
H = foot(M, B, C);
E = intersectionpoint(line(A,la), line(M,G));
F = intersectionpoint(line(B,lb), line(M,H));
D = intersectionpoint(line(C,M), line(E,F));

draw(A--B--C--cycle);
draw(E--F--M--cycle, blue);
draw(G--H);
draw(A--D--B, red);
draw(C--M);
draw(A--la);
draw(B--lb);
draw(circumcircle(triangle(E,F,G)), dashed);
draw(circumcircle(triangle(F,H,D)), dashed+purple);
draw(circumcircle(triangle(E,G,D)), dashed+purple);
draw(circumcircle(triangle(A,E,D)), dotted+red);
draw(circumcircle(triangle(B,F,D)), dotted+red);
draw(A--F);

label("$A$", A, SW);
label("$B$", B, SE);
label("$C$", C, N);
label("$M$", M, S);
label("$D$", D, NE);
label("$E$", E, W);
label("$F$", F, E);
label("$G$", G, SE);
label("$H$", H, E);
label("$\ell_1$", (A+la)/2, W);
label("$\ell_2$", (B+lb)/2, NE);
label("$\frac{64}{3}$", (E+F)/2, NE, blue);
label("$16$", (M+F)/2, SE, blue);
label("$12$", (B+D)/2, S, red);

// Clip the picture
pair lowerLeftCorner = (-9,-1); // adjust as needed
pair upperRightCorner = (9,6); // adjust as needed
clip((lowerLeftCorner.x, lowerLeftCorner.y)--(upperRightCorner.x, lowerLeftCorner.y)--(upperRightCorner.x, upperRightCorner.y)--(lowerLeftCorner.x, upperRightCorner.y)--cycle);

\end{asy}
\end{center}
Perhatikan karena $\triangle MHB \sim \triangle MBF$ dan $\triangle MGA \sim \triangle MAE$ maka $\dfrac{MG}{MA} = \dfrac{MA}{ME}$ dan $\dfrac{MH}{MB}=\dfrac{MB}{MF}$ sehingga
\begin{align*}
    MG \cdot ME = MA^2 = MB^2 = MH \cdot MF
\end{align*}
yang menunjukkan bahwa $EGHF$ siklis.  Di sisi lain, karena $\dangle CGM = \dangle CHM = 90^\circ$ maka $CGMH$ siklis juga. Oleh karena itu, kita punya
\begin{align*}
    \dangle GED = \dangle GEF = \dangle GHF = \dangle GHM = \dangle GCM = \dangle GCD
\end{align*}
sehingga $EGDC$ siklis yang langsung mengakibatkan $\dangle EDM = \dangle EDC = \dangle EGC = 90^\circ$. Ini berarti $\dangle EDM = \dangle EAM$ dan $\dangle MDF = \dangle MBF$ yang menyebabkan $EDMA$ dan $MBFD$ semuanya siklis. Dari fakta tersebut akan didapat
\begin{align*}
    \dangle MED = \dangle MAD \text{ dan } \dangle MAD = \dangle DFB = \dangle DMB
\end{align*}
yang menunjukkan $\triangle EFM \sim \triangle ADB$. Oleh karena itu
\begin{align*}
    \dfrac{AB}{EF} = \dfrac{DB}{MF} \implies AB = EF\cdot\dfrac{DB}{MF}\dfrac{64}{3}\cdot\dfrac{12}{16}= 16 \implies AM= MB = \dfrac{1}{2}AB=8
\end{align*}
dan
\begin{align*}
    BF=\sqrt{MF^2-MB^2}=\sqrt{16^2-8^2}=8\sqrt{3}
\end{align*}
sehingga
\begin{align*}
    [AMF] = \dfrac{1}{2}\cdot AM \cdot BF = \dfrac{1}{2} \cdot 8 \cdot 8\sqrt{3} = 32\sqrt{3} \implies a+b = 32 + 3 = \boxed{35}
\end{align*}
\end{proof}
